\documentclass[10pt,a4paper]{amsart}
\usepackage[margin=19mm]{geometry}
\usepackage{amsmath,amssymb,mathtools}
\usepackage[scr=rsfs,cal=euler]{mathalfa}
\usepackage{xfrac}
\usepackage[unicode,hidelinks,bookmarks=false]{hyperref}
\newcommand{\ie}{i.e.\ }
\newcommand{\cf}{cf.\ }
\newcommand{\resp}{respectively\ }
\newcommand{\Subsection}{\subsection}
\providecommand{\nobreakdash}{\nobreak\hbox{-}}
\setlength{\emergencystretch}{2em}
\multlinegap0pt
\usepackage{bm}
\usepackage{bbm}
\usepackage{dsfont}
\usepackage{xspace}
\usepackage{cancel}
\usepackage{mathtools}
\usepackage{amsthm}
\makeatletter
\@ifundefined{thm}{\newtheorem{thm}{Theorem}}{}
\@ifundefined{lem}{\newtheorem{lem}[thm]{Lemma}}{}
\@ifundefined{prop}{\newtheorem{prop}[thm]{Proposition}}{}
\makeatother
\title[Ideal Spaces of the Haagerup Tensor Product of Ternary Rings]{Ideal Spaces of the Haagerup Tensor Product of Ternary Rings of Operators}
\author{Vandana Rajpal, Arpit Kansal}
\date{Source: arXiv:2508.12367v1 (CC BY 4.0)}
\begin{document}
\maketitle

\noindent\emph{Source and licence.} Ideal Spaces of the Haagerup Tensor Product of Ternary Rings of Operators. Version arXiv:2508.12367v1 (CC BY 4.0), distributed under the Creative Commons Attribution 4.0 International licence, as recorded in the arXiv licence metadata for this exact version and shown on the abstract page https://arxiv.org/abs/2508.12367v1. Full rights notice: licenses/arxiv-2508.12367-CC-BY-4.0.txt.\par\smallskip

\noindent\textit{Editorial scope.} Counted unit: the Thm whose source label is \texttt{mainthm1} in the article (source0.tex lines 743--744, source label \texttt{mainthm1}), delivered together with the article's own complete original proof of it (source0.tex lines 746--768, one \texttt{proof} environment of 23 lines). No mathematics is written, repaired or re-derived: statement and proof are the article's own text line for line. Required context retained uncounted: fourmain (lines 695--698); result4 (lines 706--713). Every definition, display and label of those lines is retained unchanged. Omitted from this package and not redistributed: 1--694, 699--705, 714--742, 745--745, 769--1512. Nothing inside a retained excerpt is omitted or altered: every retained body is delivered exactly as the article's lines are written, so no omission receipt of any kind accompanies this package. Citations: the retained excerpts carry no citation command at all, so no bibliography entry of the article is transcribed and no reference is redistributed. Reference adaptation: none was needed and none was made; every reference inside the delivered unit resolves inside it, so no unresolved reference or citation is produced. Printed slips kept literal: no printed word, symbol, hypothesis or number of the counted statement or of its proof was altered. Layout and class: the article's own class is \texttt{birkjour} with packages including tikz-cd, amsmath, amssymb, url, graphicx, xy; the wrapper is the lane's pinned \texttt{amsart} editorial wrapper at 10pt on A4, which restates the theorem-like environments the retained text needs and adds the article's own macro definitions that the retained text needs (none), with extra wrapper packages (bm, bbm, dsfont, xspace, cancel, mathtools). The delivered numbering is the unit's own. Assets: no retained span contains a figure environment, a graphics-inclusion command, a PDF-page inclusion, a table or a TikZ picture; no image file is copied into this package, the package declares no asset dependency and no image file is redistributed with it. Licence: version arXiv:2508.12367v1 of this article is distributed under the Creative Commons Attribution 4.0 International licence, as recorded in the arXiv licence metadata for this exact version and shown on the abstract page https://arxiv.org/abs/2508.12367v1. A full rights notice is delivered at licenses/arxiv-2508.12367-CC-BY-4.0.txt. Characters: every retained excerpt is pure ASCII, so the delivered engine, XeLaTeX with its default Latin Modern text font, reports no missing character.
\begin{prop} \label{fourmain}
    The restriction of $\Phi$ to $\text{Prim}(V) \times \text{Prim}(B)$ is continuous for the 
product \textit{hk}-topology on $\text{Prim}(V) \times \text{Prim}(B)$ and the $\tau_w$-topology on $\text{Prim}(V \otimes^h B)$.
\end{prop}

\begin{lem}\label{result4}

Let $I_0$ be a proper ideal of $V$ and $J_0$ a proper ideal of $B$. Then, 

\[
V\otimes^{h} J_0 + I_0 \otimes^{h} B = \bigcap \{ V\otimes^{h} Q + P \otimes^{h} B : P \in hull(I_0), Q\in hull(J_0) \}.
\]
\end{lem}

\begin{thm} \label{mainthm1} The mapping $\Phi: \operatorname{Id}(V) \times \operatorname{Id}(B) \to \operatorname{Id}(V \otimes^h B)$ is $\tau_w$-continuous.
\end{thm}

\begin{proof}
      We need  to show that if $K \in \operatorname{Id}(V \otimes^h B)$ and $I_0 \in \operatorname{Id}(V)$, $J_0 \in \operatorname{Id}(B)$ satisfy $\Phi(I_0, J_0) \in U(K)$,  there exist $\tau_w$-neighbourhoods $V$ of $I_0$ and $W$ of $J_0$ in $\operatorname{Id}(V)$ and $\operatorname{Id}(B)$, respectively, such that $\Phi(V \times W) \subseteq U(K)$. 

 Since $I_0$ and $J_0$ are necessarily proper, it follows from  Lemma \ref{result4} that $V\otimes^h J_0 + I_0 \otimes^h B$ is the intersection of ideals of the form
\[
V \otimes^h Q + P \otimes^h B,
\]
where $P \in h(I_0)$ and $Q \in h(J_0)$. Thus, since $K \nsubseteq V \otimes^h J_0+ I_0 \otimes^h B$,  there exists $P_0 \in h(I_0)$ and $Q_0 \in h(J_0)$ such that
 $K \nsubseteq (V \otimes^h Q_0 + P_0 \otimes^h B)$.
Since $\Phi$ is continuous on $\text{Prim}(V) \times \text{Prim}(B)$ (Proposition \ref{fourmain}) and
\[
\Phi(P_0, Q_0) \in U(K) \cap \text{Prim}(V \otimes^h B),
\]
there exist open neighbourhoods $V_0$ of $P_0$ in $\text{Prim}(V)$ and $W_0$ of $Q_0$ in $\text{Prim}(B)$ such that $\Phi(V_0 \times W_0) \subseteq U(K)$. Consequently, $K \nsubseteq V \otimes^h Q + P \otimes^h B$ for all $P \in V_0$ and $Q \in W_0$. Now, define
\[
U_1 = \{I \in \operatorname{Id}(V): h(I) \cap V_0 \neq \emptyset\}, \quad U_2 = \{J \in \operatorname{Id}(B): h(J) \cap W_0 \neq \emptyset\}.
\]
Then $I_0 \in U_1$ and $J_0 \in U_2$, and $U_1$,  $U_2$ are $\tau_w$-open in $\operatorname{Id}(V)$ and  $\operatorname{Id}(B)$, respectively. Finally, let $I \in U_1$ and $J \in U_2$. Then, there exist $P \in h(I) \cap V_0$ and $Q \in h(J) \cap W_0$. Then
\[
V \otimes^h J + I \otimes^h B \subseteq V \otimes^h Q + P \otimes^h B \quad \text{and} \quad  K \nsubseteq V \otimes^h Q + P \otimes^h B.
\]
This establishes that $V \otimes^h J + I \otimes^h B \in U(K)$, completing the proof. 
\end{proof}

\end{document}
